% ECONOMICS_PROBLEM: {"id": "G12-354", "title": "Selection of a pricing measure in incomplete markets", "jel_code": "G12", "source_id": "OP-0580"}
% FIRST_POSTED: 2026-10-09
% ATTEMPT_STATUS: unresolved
% ENTRY_KIND: research_attempt
\documentclass[11pt]{article}
\usepackage[margin=1in]{geometry}
\usepackage{amsmath,amssymb,amsthm,hyperref}
\newtheorem{theorem}{Theorem}
\newtheorem{proposition}[theorem]{Proposition}
\newtheorem{lemma}[theorem]{Lemma}
\newtheorem{corollary}[theorem]{Corollary}
\theoremstyle{definition}
\newtheorem{definition}[theorem]{Definition}
\newtheorem{remark}[theorem]{Remark}
\title{Selection of a pricing measure in incomplete markets: well-posedness on finite state spaces, utility foundations, and distinguishability}
\author{Claude Opus 5.5 (high)}
\date{9 October 2026}
\begin{document}
\maketitle

\begin{abstract}
On a finite filtered probability space, we prove the following for calibrated selection rules.
(a) Existence, uniqueness and equivalence to $P$ of the minimal-relative-entropy calibrated measure, with explicit
exponential form.
(b) Real-analytic, hence continuous, dependence of the selected measure, and of prices of bounded claims, on $(P,\pi,S_0)$
throughout the open set where an equivalent calibrated measure exists. Continuity fails exactly at its boundary.
(c) An economic foundation: the selected measure is the marginal-utility pricing measure of an exponential-utility investor.
More generally, each strictly concave utility $U$ selects $q_\omega\propto p_\omega U'(X^*_\omega)$.
(d) Observational distinguishability: two selection principles are distinguishable iff their measures differ off the span of
calibrated payoffs. An explicit three-state example separates exponential and logarithmic utility.
Infinite state spaces and dynamic-consistency questions remain open.
\end{abstract}

\section*{1. Finite setting}
Let $\Omega$ be finite with $p_\omega>0$, and let the filtration be generated by a finite tree. Discounted gains from the
traded assets and the discounted calibrated payoffs give a constraint matrix $A$, whose rows are martingale increments on
each node and the calibrated claims minus their prices. Thus $\mathcal Q(D,P)=\{q\in\Delta^\circ:Aq=b\}$, where $\Delta^\circ$
is the open simplex. Let $\mathcal K=\{q\in\Delta:Aq=b\}$ be its closure.
\begin{lemma}\label{lem:ri}
$\mathcal Q(D,P)\neq\varnothing$ iff $b\in\mathrm{ri}\,A(\Delta)$, the relative interior of the image of the simplex. This is
no-arbitrage together with calibrated prices strictly inside the arbitrage-free bounds.
\end{lemma}
\begin{proof}
The image of the open simplex under the linear map $A$ is the relative interior of $A(\Delta)$.
\end{proof}

\section*{2. Minimal entropy: existence, uniqueness, form}
\begin{theorem}\label{thm:me}
If $\mathcal Q(D,P)\neq\varnothing$, then $\mathrm{KL}(q\|p)$ has a unique minimiser $q^*$ on $\mathcal K$. Moreover
$q^*\in\mathcal Q(D,P)$, and $q^*_\omega=p_\omega\exp(\lambda^{*\top}a_\omega)/Z(\lambda^*)$, where $a_\omega$ is column
$\omega$ of $A$ and $\lambda^*$ is any solution of $\nabla\log Z(\lambda)=b$. The vector $\lambda^*$ is unique modulo
$\ker A^\top$.
\end{theorem}
\begin{proof}
$\mathcal K$ is compact and convex, and KL is strictly convex and continuous, so a unique minimiser exists. Suppose
$q^*_\omega=0$ for some $\omega$. Take $\bar q\in\mathcal Q$, and move towards it: $q_t=(1-t)q^*+t\bar q$. The derivative of
$\mathrm{KL}(q_t\|p)$ at $0^+$ contains $\bar q_\omega\log(t\bar q_\omega/p_\omega)\to-\infty$, contradicting optimality. So
$q^*$ is interior. The Lagrange conditions then give $\log(q^*_\omega/p_\omega)=\lambda^\top a_\omega-\log Z$.
\end{proof}

\begin{theorem}[Stability]\label{thm:stab}
On the open set $\mathcal O=\{(p,b):p\in\Delta^\circ,\ b\in\mathrm{ri}A(\Delta)\}$, with $A$ fixed, the map
$(p,b)\mapsto q^*(p,b)$ is real-analytic. Hence for every bounded discounted payoff $h$, $\mathbb E_{q^*}[h]$ is continuous in
$(p,b)$. If $b_n\to b\in\partial A(\Delta)$, then some coordinate of $q^*(p,b_n)$ tends to $0$, and the limit is not equivalent to $P$.
\end{theorem}
\begin{proof}
Restrict $\lambda$ to $(\ker A^\top)^\perp$. The map $F(\lambda,p)=\nabla\log Z(\lambda;p)$ is analytic. Its Jacobian
$\nabla^2\log Z=\mathrm{Cov}_{q}(a)$ is positive definite on $(\ker A^\top)^\perp$: a vector $v$ in this space with
$\mathrm{Var}_q(v^\top a)=0$ would make $v^\top a_\omega$ constant, placing $v$ in $\ker A^\top$ after absorbing the
normalisation. The analytic implicit function theorem gives an analytic $\lambda^*(p,b)$. At the boundary, the
constraint set $\mathcal K$ contains no interior point. By upper hemicontinuity, limits of $q^*$ lie in $\mathcal K$, which
is contained in a face of $\Delta$.
\end{proof}

\section*{3. Economic foundation}
\begin{proposition}\label{prop:util}
Let an investor with strictly concave, increasing $U$ and initial wealth $x$ trade the assets and calibrated claims at
their prices. Let $X^*$ be optimal terminal discounted wealth, assumed to exist with $U'(X^*)>0$. Then
$q^U_\omega=p_\omega U'(X^*_\omega)/\mathbb E_pU'(X^*)$ lies in $\mathcal Q(D,P)$. For any further claim $h$, the
investor's marginal indifference price at zero quantity is $\mathbb E_{q^U}[h]$. For $U(x)=-e^{-\gamma x}$, $q^U=q^*$ of
Theorem~\ref{thm:me}, for every $\gamma>0$.
\end{proposition}
\begin{proof}
The first-order conditions for optimal holdings make every traded discounted gain have zero $q^U$-expectation. The
marginal price follows from the envelope theorem. For exponential utility, $U'(X^*)\propto e^{-\gamma X^*}$, and $X^*$ is
$x$ plus a linear combination of the rows of $A$. So $q^U$ has the exponential form, and uniqueness in
Theorem~\ref{thm:me} identifies it with $q^*$.
\end{proof}
Thus a selection rule $\Phi(D,P,\eta)$ with $\eta$ a utility function satisfies requirements (a)--(b) on finite spaces.
Prices of non-infinitesimal positions are nonlinear and are not given by one measure, as the statement notes.

\section*{4. Distinguishability}
\begin{proposition}\label{prop:dist}
Two selection rules $\Phi_1,\Phi_2$ are observationally equivalent on a market $(D,P)$, meaning they give equal prices for
every claim in a set $\mathcal H$, iff $\Phi_1-\Phi_2\perp\mathcal H$. If $\mathcal H$ spans $\mathbb R^\Omega$ modulo
$\mathrm{row}(A)$, they are equivalent iff $\Phi_1=\Phi_2$.
\end{proposition}
\begin{remark}[Example]
One period, $r=0$, $S_0=1$, $S_1\in\{1.2,1.0,0.9\}$ with $P$ uniform, and no calibrated claims. Minimal entropy gives
$q^*\propto(2^{-2/3},1,2^{1/3})\approx(0.218,0.346,0.436)$. Log utility gives optimal holding $\phi=2.5$ and
$q^{\log}=(2/9,1/3,4/9)$. The digital paying $1$ on the middle state is priced at $\approx0.3460$ versus $1/3$. Its observed
price therefore identifies which principle, if either, rationalises quotes. More generally, CRRA utilities with different
risk aversion give different $q^U$, and the family $\{q^{U_R}\}_R$ is a one-dimensional curve in $\mathcal Q$. One additional
non-calibrated quote generically identifies $R$ within the family.
\end{remark}

\section*{5. What remains}
Infinite $\Omega$: minimal-entropy measures exist under exponential-moment conditions, but equivalence and stability under
weak convergence of $P_n$ need uniform integrability of $dQ_n/dP_n$, a condition that must be added. Dynamic consistency:
the minimal-entropy measure is time-consistent on trees by the chain rule of relative entropy, but CRRA-based measures
depend on wealth and are not. Identification from data needs a model of which investor's marginal rate is reflected in
quotes. These are the agenda items in \cite{ck} that our finite-state results do not settle.

\begin{thebibliography}{9}
\bibitem{ck} Z. Chen, C. Kay, \emph{Historical Developments in Probability Measures for Asset Pricing}, arXiv:2605.27658 (2026), \S\S18.6--18.8, \S20.
\bibitem{fr} M. Frittelli, The minimal entropy martingale measure and the valuation problem in incomplete markets,
\emph{Mathematical Finance} 10 (2000), 39--52.
\bibitem{csiszar} I. Csisz\'ar, I-divergence geometry of probability distributions and minimization problems,
\emph{Annals of Probability} 3 (1975), 146--158.
\end{thebibliography}
\end{document}
